% ============================================================
\section{Related Work}
% ============================================================
\label{sec:related}

\subsection{Privacy Inference Attacks on Language Models}

The capability of language models to infer personal attributes from unstructured text builds on a longer tradition of computational author profiling research~\citep{koppel2002automatically,rangel2013overview,estival2007author}.
\citet{staab2024beyond} conducted the first systematic evaluation of LLMs as privacy-violating inference engines, demonstrating that models such as GPT-4 can infer age, gender, location, and other sensitive attributes from Reddit-style posts with accuracy substantially exceeding random baselines. Their work establishes the SynthPAI corpus and the hardness taxonomy that underpin our evaluation. Crucially, however, \citet{staab2024beyond} evaluate inference \emph{capability}---what LLMs can do when given a profiling prompt---rather than \emph{detection}---whether such profiling prompts can be identified and blocked before reaching the model. Our work addresses this orthogonal and previously unstudied question.

A related line of work probes LLMs for leakage of personally identifiable information rather than open-ended attribute inference: \citet{lukas2023analyzing} analyze how much PII language models leak and to what extent it can be mitigated, and \citet{kim2023propile} introduce ProPILE, a tool that probes whether a target LLM will surface a specific individual's known PII when prompted. Both target information the model may have memorized about real individuals during training, whereas UAL-Inference targets attributes of an arbitrary author never seen during training, synthesized purely from the text supplied at inference time.

A parallel line of research on membership inference and training data extraction~\cite{carlini2021extracting} targets the recovery of verbatim training data through repeated querying. These attacks exploit memorization rather than reasoning, and are therefore conceptually and operationally distinct from UAL-Inference, which requires no memorization and succeeds on information entirely absent from training data.

\subsection{Prompt Injection and Jailbreak Defenses}

Prompt injection attacks exploit the inability of LLMs to distinguish between instruction and data in the input stream~\cite{perez2022ignore,hui2024pleak}. Several input-boundary defenses have been proposed to detect such attacks~\cite{agarwal2024prompt}. Prompt Guard~2~\cite{metaai2024promptguard2} and Llama Guard~\cite{inan2023llama} are the most widely deployed classifier-based mechanisms, using DeBERTa-v3 and Llama fine-tuned on labeled injection and jailbreak datasets. These classifiers achieve strong detection performance on the threat class they were designed for.

Our work extends this line by demonstrating a \emph{systematic blind spot}: the feature space of injection classifiers does not generalize to UAL-Inference, a semantically distinct threat category that carries no syntactic injection markers. This observation motivates the shift from syntax-based to intent-based detection, which is the core contribution of UAL Semantic Guard.

\subsection{LLM-as-a-Judge Frameworks}

Using LLMs as zero-shot evaluators---a paradigm known as LLM-as-a-Judge---has been studied primarily in the context of answer quality assessment and alignment evaluation~\cite{zheng2023judging}. Recent work has extended this paradigm to safety evaluation~\cite{dong2024safeguarding}, where a secondary LLM audits the output of a primary model for policy violations. Our use of LLM-as-a-Judge for \emph{input}-side intent classification---rather than output auditing---is, to the best of our knowledge, novel in the privacy defense literature. We further demonstrate that zero-shot prompting of a local 1.5B-parameter model suffices to achieve 100\% UAL detection, suggesting that the task is well-suited to the few-shot in-context learning capabilities of current open-weight models.

\subsection{Anonymization as a Privacy Defense}

The ETH-SRI benchmark~\cite{ethsri-llmprivacy-repo} evaluates named-entity anonymization (replacing explicit personal references with typed placeholders) as a defense against attribute inference. This approach targets \emph{explicit} attribute mentions---reducing, for example, the availability of location names in the text---but does not prevent the inference of attributes that are implicit or encoded in writing style. Furthermore, anonymization operates on the input \emph{document} rather than on the profiling \emph{prompt}, and therefore does not constitute an input-boundary defense in the sense evaluated in our work. We treat it as background context motivating the need for prompt-level intent detection.
